\renewcommand{\currentchapter}{8}
\title{Analisis Kesenjangan Pembimbingan}
\subtitle{Bab 8 \textbar\ Dari kondisi 2026 menuju sistem belajar yang terukur pada 2027}
\date{SMAN 1 Muntilan bersama StudentPro \textbar\ Tahun Ajaran 2026/2027}

\begingroup
\setbeamercolor{background canvas}{bg=StudentNavy}
\begin{frame}[plain]
  \vspace{0.65cm}
  \kicker{SMAN 1 Muntilan \textbar\ StudentPro}
  \vspace{0.08cm}
  {\usebeamerfont{title}\color{StudentWhite}\inserttitle\par}
  \vspace{0.38cm}
  {\usebeamerfont{subtitle}\color{StudentSoft}\insertsubtitle\par}
  \vfill
  {\color{StudentGold}\rule{2.3cm}{2.5pt}}\par
  \vspace{0.20cm}
  {\usebeamerfont{author}\color{StudentWhite}\insertauthor}\par
  {\usebeamerfont{date}\color{StudentSoft}\insertdate}
  \vspace{0.42cm}
\end{frame}
\endgroup

\begin{frame}{Bab 8 mengubah evaluasi menjadi agenda perbaikan}
  \kicker{Posisi dalam Laporan}
  \begin{columns}[T,onlytextwidth]
    \begin{column}{0.30\textwidth}
      \begin{block}{Bab 5}
        Menilai capaian, bukti, keterbatasan, dan pelajaran program 2026.
      \end{block}
    \end{column}
    \begin{column}{0.30\textwidth}
      \begin{block}{Bab 8}
        Menentukan jarak antara kondisi saat ini dan sistem yang dibutuhkan.
      \end{block}
    \end{column}
    \begin{column}{0.30\textwidth}
      \begin{block}{Bab 9}
        Menetapkan target 2027 yang terukur dan dapat dipantau.
      \end{block}
    \end{column}
  \end{columns}
  \vspace{0.26cm}
  \claim{Kesenjangan tidak digunakan untuk menyalahkan, tetapi untuk memilih perbaikan yang paling berdampak.}
\end{frame}

\begin{frame}{Kesenjangan memerlukan dua titik yang jelas}
  \kicker{Definisi Operasional}
  \begin{columns}[T,onlytextwidth]
    \begin{column}{0.44\textwidth}
      \begin{block}{Kondisi saat ini}
        Bukti yang tersedia tentang kemampuan siswa, pelaksanaan, tindak lanjut, data, dan hasil.
      \end{block}
    \end{column}
    \begin{column}{0.44\textwidth}
      \begin{block}{Kondisi yang dibutuhkan}
        Standar internal agar setiap siswa memperoleh pemetaan, latihan, remedial, konsultasi, dan pelaporan yang tepat.
      \end{block}
    \end{column}
  \end{columns}
  \vspace{0.20cm}
  \begin{beamercolorbox}[rounded=true,sep=0.20cm,wd=\textwidth]{block body}
    \centering\bfseries Kesenjangan = kebutuhan yang belum terpenuhi dan dibuktikan dengan data yang dapat ditelusuri.
  \end{beamercolorbox}
\end{frame}

\begin{frame}[shrink=6]{Status bukti menentukan kekuatan simpulan}
  \kicker{Batas Analisis}
  {\fontsize{10.0}{12.5}\selectfont
  \renewcommand{\arraystretch}{0.98}
  \begin{tabularx}{\textwidth}{>{\bfseries\color{StudentBlue}}p{0.23\textwidth} >{\RaggedRight\arraybackslash}X >{\RaggedRight\arraybackslash}p{0.29\textwidth}}
    \toprule
    \textbf{Status} & \textbf{Bukti} & \textbf{Perlakuan} \\
    \midrule
    Terverifikasi & 124 siswa unik tercatat menerima konsultasi akademik & Dapat dinyatakan sebagai capaian proses \\
    Tersedia internal & Rekap hasil seleksi, aktivitas, dan catatan ranking & Digunakan terbatas setelah validasi identitas dan metadata \\
    Belum lengkap & Kohor final, baseline seragam, tren tujuh subtes, dan status remedial & Tidak dibuat persentase atau kesimpulan kuantitatif \\
    Acuan resmi & Framework UTBK--SNBT 2026 & Menetapkan kemampuan yang perlu dipetakan \citep{snpmb_framework_2026} \\
    \bottomrule
  \end{tabularx}}
  \vspace{0.10cm}
  {\color{StudentGray}\fontsize{9.5}{11.8}\selectfont Sumber internal dirinci pada bibliografi di bagian paling akhir dokumen.}
\end{frame}

\begin{frame}{Enam kesenjangan dibaca sebagai satu sistem}
  \kicker{Peta Utama}
  \begin{columns}[T,onlytextwidth]
    \begin{column}{0.30\textwidth}
      \begin{block}{1. Data}
        Identitas, baseline, aktivitas, dan hasil belum sepenuhnya terhubung.
      \end{block}
      \begin{block}{2. Akademik}
        Penguasaan tujuh subtes belum tergambar seragam.
      \end{block}
    \end{column}
    \begin{column}{0.30\textwidth}
      \begin{block}{3. Proses}
        Ketuntasan materi dan latihan belum selalu dapat ditelusuri.
      \end{block}
      \begin{block}{4. Tindak lanjut}
        Kesalahan belum selalu berakhir pada remedial dan tes ulang.
      \end{block}
    \end{column}
    \begin{column}{0.30\textwidth}
      \begin{block}{5. Strategi}
        Waktu, ketelitian, dan stamina perlu dibaca bersama skor.
      \end{block}
      \begin{block}{6. Pilihan studi}
        Target prodi perlu dihubungkan dengan profil dan peluang siswa.
      \end{block}
    \end{column}
  \end{columns}
\end{frame}

\begin{frame}{Kesenjangan data dimulai dari daftar peserta}
  \kicker{Gap 1 -- Satu Kohor yang Valid}
  \begin{columns}[T,onlytextwidth]
    \begin{column}{0.46\textwidth}
      {\color{StudentBlue}\bfseries Kondisi yang ditemukan}\par
      \begin{itemize}
        \item Penerima konsultasi tidak sama dengan seluruh peserta program.
        \item Catatan kegiatan dan hasil berada pada beberapa rekap.
        \item Jalur SNBT dan mandiri berisiko tercampur.
      \end{itemize}
    \end{column}
    \begin{column}{0.48\textwidth}
      {\color{StudentBlue}\bfseries Kondisi yang dibutuhkan}\par
      \begin{itemize}
        \item Identitas peserta unik dan status kohor.
        \item Riwayat kegiatan dalam satu profil.
        \item Hasil seleksi dipisahkan per jalur.
      \end{itemize}
    \end{column}
  \end{columns}
  \vspace{0.16cm}
  \claim{Tanpa kohor final, target 70\% belum dapat dibandingkan dengan capaian secara sah.}
\end{frame}

\begin{frame}{Baseline harus tersedia untuk setiap siswa dan subtes}
  \kicker{Gap 1 -- Titik Awal}
  \begin{columns}[T,onlytextwidth]
    \begin{column}{0.30\textwidth}
      \begin{block}{Cakupan}
        Apakah seluruh siswa sudah mengikuti diagnostik tujuh subtes?
      \end{block}
    \end{column}
    \begin{column}{0.30\textwidth}
      \begin{block}{Keterbandingan}
        Apakah instrumen, durasi, dan indikator dicatat dengan cara yang sama?
      \end{block}
    \end{column}
    \begin{column}{0.30\textwidth}
      \begin{block}{Keterlacakan}
        Apakah hasil awal terhubung dengan target dan program remedial?
      \end{block}
    \end{column}
  \end{columns}
  \vspace{0.10cm}
  \begin{beamercolorbox}[rounded=true,sep=0.16cm,wd=\textwidth]{block body}
    \centering\bfseries Nilai terbaru tanpa baseline menunjukkan posisi, tetapi belum menunjukkan kemajuan.
  \end{beamercolorbox}
\end{frame}

\begin{frame}[shrink=8]{Tujuh subtes membutuhkan tujuh peta kesenjangan}
  \kicker{Gap 2 -- Pemetaan Akademik}
  {\fontsize{9.25}{11.5}\selectfont
  \renewcommand{\arraystretch}{0.93}
  \begin{tabularx}{\textwidth}{>{\bfseries\color{StudentBlue}}p{0.25\textwidth} >{\RaggedRight\arraybackslash}X >{\RaggedRight\arraybackslash}p{0.30\textwidth}}
    \toprule
    \textbf{Subtes} & \textbf{Bukti minimal} & \textbf{Pertanyaan kesenjangan} \\
    \midrule
    Penalaran Umum & Akurasi induktif, deduktif, kuantitatif & Cara berpikir mana yang belum stabil? \\
    PPU & Makna, hubungan gagasan, pengetahuan konteks & Apakah bahasa atau pengetahuan konteks menghambat? \\
    PBM & Kalimat efektif, kepaduan, organisasi teks & Kesalahan terjadi pada kaidah atau struktur gagasan? \\
    Pengetahuan Kuantitatif & Bilangan, aljabar, geometri, data & Konsep atau prosedur mana yang lambat? \\
    LBI & Proses literasi dan konteks bacaan & Konteks mana yang paling menahan hasil? \\
    LBE & Informasi, inferensi, argumen, kosakata & Apakah masalah bahasa atau strategi membaca? \\
    Penalaran Matematika & Formulasi, penggunaan, interpretasi & Tahap pemodelan mana yang gagal? \\
    \bottomrule
  \end{tabularx}}
\end{frame}

\begin{frame}{Skor subtes belum cukup untuk menentukan intervensi}
  \kicker{Tiga Dimensi Penguasaan}
  \begin{columns}[T,onlytextwidth]
    \begin{column}{0.30\textwidth}
      \begin{block}{Akurasi}
        Seberapa sering siswa mengambil keputusan yang benar?
      \end{block}
    \end{column}
    \begin{column}{0.30\textwidth}
      \begin{block}{Waktu}
        Apakah proses benar dapat dilakukan dalam batas yang tersedia?
      \end{block}
    \end{column}
    \begin{column}{0.30\textwidth}
      \begin{block}{Stabilitas}
        Apakah hasil bertahan pada beberapa tes dan simulasi penuh?
      \end{block}
    \end{column}
  \end{columns}
  \vspace{0.26cm}
  \claim{Kesenjangan yang sama nilainya dapat membutuhkan tindakan yang berbeda.}
\end{frame}

\begin{frame}{Empat profil menunjukkan kebutuhan yang berbeda}
  \kicker{Membaca Akurasi dan Waktu}
  {\fontsize{10.1}{12.6}\selectfont
  \renewcommand{\arraystretch}{0.98}
  \begin{tabularx}{\textwidth}{>{\bfseries\color{StudentBlue}\RaggedRight\arraybackslash}p{0.23\textwidth} >{\RaggedRight\arraybackslash}X >{\RaggedRight\arraybackslash}p{0.30\textwidth}}
    \toprule
    \textbf{Profil} & \textbf{Interpretasi awal} & \textbf{Tindak lanjut} \\
    \midrule
    Akurat dan cepat & Penguasaan relatif stabil & Pengayaan dan integrasi subtes \\
    Akurat tetapi lambat & Strategi belum efisien & Latihan waktu dan perbandingan cara \\
    Cepat tetapi tidak akurat & Keputusan terlalu agresif & Verifikasi, kontrol tempo, dan ketelitian \\
    Lambat dan tidak akurat & Konsep atau pemahaman stimulus lemah & Materi dasar dan latihan bertahap \\
    \bottomrule
  \end{tabularx}}
\end{frame}

\begin{frame}{Kesenjangan literasi perlu dipisah menurut proses}
  \kicker{Gap 2 -- Bahasa dan Bacaan}
  \begin{columns}[T,onlytextwidth]
    \begin{column}{0.46\textwidth}
      {\color{StudentBlue}\bfseries Gejala yang terlihat}\par
      \begin{itemize}
        \item Kehilangan gagasan utama.
        \item Salah memahami rujukan atau makna kontekstual.
        \item Memilih klaim yang terlalu mutlak.
        \item Tidak membedakan fakta dan inferensi.
      \end{itemize}
    \end{column}
    \begin{column}{0.48\textwidth}
      {\color{StudentBlue}\bfseries Data yang diperlukan}\par
      \begin{itemize}
        \item Indikator soal dan jenis proses literasi.
        \item Konteks bacaan LBI.
        \item Bahasa, panjang, dan kompleksitas teks.
        \item Bukti pilihan salah yang berulang.
      \end{itemize}
    \end{column}
  \end{columns}
  \vspace{0.14cm}
  \begin{beamercolorbox}[rounded=true,sep=0.17cm,wd=\textwidth]{block body}
    \centering\bfseries ``Lemah literasi'' harus diuraikan menjadi proses dan konteks yang dapat dilatih.
  \end{beamercolorbox}
\end{frame}

\begin{frame}{Kesenjangan numerasi perlu dipisah menurut tahap berpikir}
  \kicker{Gap 2 -- Kuantitatif dan Matematika}
  \begin{columns}[T,onlytextwidth]
    \begin{column}{0.46\textwidth}
      {\color{StudentBlue}\bfseries Gejala yang terlihat}\par
      \begin{itemize}
        \item Salah memilih konsep atau rumus.
        \item Model tidak sesuai konteks.
        \item Prosedur benar tetapi terlalu panjang.
        \item Hasil tidak diperiksa kewajarannya.
      \end{itemize}
    \end{column}
    \begin{column}{0.48\textwidth}
      {\color{StudentBlue}\bfseries Data yang diperlukan}\par
      \begin{itemize}
        \item Domain materi dan prasyarat.
        \item Tahap formulasi, penggunaan, interpretasi.
        \item Waktu per soal dan langkah penyelesaian.
        \item Jenis kesalahan hitung atau model.
      \end{itemize}
    \end{column}
  \end{columns}
  \vspace{0.14cm}
  \claim{Nilai rendah tidak selalu berarti kekurangan rumus; masalah dapat muncul sebelum perhitungan dimulai.}
\end{frame}

\begin{frame}{Aktivitas belajar belum otomatis menjadi penguasaan}
  \kicker{Gap 3 -- Proses Pembelajaran}
  \begin{columns}[T,onlytextwidth]
    \begin{column}{0.46\textwidth}
      \begin{block}{Yang mudah dicatat}
        Kehadiran, materi dibuka, kuis dikerjakan, jumlah soal, dan skor akhir.
      \end{block}
    \end{column}
    \begin{column}{0.48\textwidth}
      \begin{block}{Yang harus dibuktikan}
        Konsep dipahami, kesalahan berkurang, strategi membaik, dan hasil stabil pada tes baru.
      \end{block}
    \end{column}
  \end{columns}
  \vspace{0.22cm}
  \begin{beamercolorbox}[rounded=true,sep=0.20cm,wd=\textwidth]{block body}
    \centering\bfseries Ketuntasan aktivitas adalah syarat proses; ketuntasan kemampuan memerlukan bukti tes.
  \end{beamercolorbox}
\end{frame}

\begin{frame}{Remedial yang terlambat memperlebar kesenjangan}
  \kicker{Gap 4 -- Tindak Lanjut}
  {\fontsize{10.1}{12.6}\selectfont
  \renewcommand{\arraystretch}{0.99}
  \begin{tabularx}{\textwidth}{>{\bfseries\color{StudentBlue}}p{0.18\textwidth} >{\RaggedRight\arraybackslash}X >{\RaggedRight\arraybackslash}p{0.29\textwidth}}
    \toprule
    \textbf{Tahap} & \textbf{Kondisi yang dibutuhkan} & \textbf{Bukti selesai} \\
    \midrule
    Deteksi & Kesalahan teridentifikasi segera setelah tes & Label indikator dan penyebab \\
    Penugasan & Materi serta latihan sesuai penyebab & Tugas memiliki tenggat dan penanggung jawab \\
    Pelaksanaan & Siswa menyelesaikan bantuan terarah & Aktivitas remedial tercatat \\
    Tes ulang & Kemampuan diuji dengan soal setara & Perubahan akurasi dan waktu \\
    Keputusan & Hasil menentukan lanjut atau ulang & Status tuntas, remedial, atau pengayaan \\
    \bottomrule
  \end{tabularx}}
\end{frame}

\begin{frame}{Monitoring lemah ketika laporan tidak menghasilkan keputusan}
  \kicker{Gap 4 -- Siklus yang Belum Tertutup}
  \begin{columns}[T,onlytextwidth]
    \begin{column}{0.30\textwidth}
      \begin{block}{Data ada}
        Skor dan aktivitas tercatat, tetapi belum diprioritaskan.
      \end{block}
    \end{column}
    \begin{column}{0.30\textwidth}
      \begin{block}{Tugas ada}
        Remedial diberikan, tetapi status selesai belum diperiksa.
      \end{block}
    \end{column}
    \begin{column}{0.30\textwidth}
      \begin{block}{Hasil ada}
        Tes ulang tersedia, tetapi keputusan berikutnya belum dicatat.
      \end{block}
    \end{column}
  \end{columns}
  \vspace{0.26cm}
  \claim{Monitoring dinyatakan selesai setelah tindakan diperiksa kembali, bukan setelah laporan dikirim.}
\end{frame}

\begin{frame}{Strategi waktu dan stamina perlu menjadi data}
  \kicker{Gap 5 -- Kinerja Ujian}
  \begin{columns}[T,onlytextwidth]
    \begin{column}{0.46\textwidth}
      {\color{StudentBlue}\bfseries Indikator strategi}\par
      \begin{itemize}
        \item Waktu per soal dan subtes.
        \item Soal kosong dan tebakan.
        \item Perubahan tempo selama tes.
        \item Akurasi sebelum dan sesudah titik lelah.
      \end{itemize}
    \end{column}
    \begin{column}{0.48\textwidth}
      {\color{StudentBlue}\bfseries Keputusan latihan}\par
      \begin{itemize}
        \item Latihan topikal atau paket waktu.
        \item Urutan pengerjaan internal.
        \item Teknik melewati hambatan.
        \item Simulasi bertahap menuju tes penuh.
      \end{itemize}
    \end{column}
  \end{columns}
  \vspace{0.15cm}
  \claim{Kesenjangan stamina tidak terlihat pada kuis singkat yang berdiri sendiri.}
\end{frame}

\begin{frame}{Pilihan program studi harus terhubung dengan profil siswa}
  \kicker{Gap 6 -- Arah Studi}
  \begin{columns}[T,onlytextwidth]
    \begin{column}{0.30\textwidth}
      \begin{block}{Aspirasi}
        Minat, program studi, kampus, dan alasan pilihan siswa.
      \end{block}
    \end{column}
    \begin{column}{0.30\textwidth}
      \begin{block}{Kesiapan}
        Profil tujuh subtes, konsistensi, rekam belajar, dan kemampuan pendukung.
      \end{block}
    \end{column}
    \begin{column}{0.30\textwidth}
      \begin{block}{Strategi}
        Pilihan utama, alternatif rasional, jalur seleksi, serta jadwal keputusan.
      \end{block}
    \end{column}
  \end{columns}
  \vspace{0.26cm}
  \begin{beamercolorbox}[rounded=true,sep=0.18cm,wd=\textwidth]{block body}
    \centering\bfseries Konsultasi menghubungkan cita-cita dengan tindakan akademik yang konkret.
  \end{beamercolorbox}
\end{frame}

\begin{frame}{Prioritas ditentukan oleh dampak, urgensi, dan kendali}
  \kicker{Matriks Keputusan}
  \begin{columns}[T,onlytextwidth]
    \begin{column}{0.30\textwidth}
      \begin{block}{Dampak}
        Seberapa besar kesenjangan menahan kemajuan atau keputusan siswa?
      \end{block}
    \end{column}
    \begin{column}{0.30\textwidth}
      \begin{block}{Urgensi}
        Seberapa dekat tenggat belajar, tryout, pemilihan prodi, atau pendaftaran?
      \end{block}
    \end{column}
    \begin{column}{0.30\textwidth}
      \begin{block}{Kendali}
        Apakah sekolah, tentor, siswa, dan keluarga dapat melakukan perubahan sekarang?
      \end{block}
    \end{column}
  \end{columns}
  \vspace{0.26cm}
  \claim{Kesenjangan terbesar belum tentu dikerjakan pertama jika bukti atau prasyaratnya belum tersedia.}
\end{frame}

\begin{frame}{Tiga tingkat prioritas menjaga fokus program}
  \kicker{Klasifikasi Tindak Lanjut}
  {\fontsize{10.2}{12.7}\selectfont
  \renewcommand{\arraystretch}{1.00}
  \begin{tabularx}{\textwidth}{>{\bfseries\color{StudentBlue}}p{0.18\textwidth} >{\RaggedRight\arraybackslash}X >{\RaggedRight\arraybackslash}p{0.30\textwidth}}
    \toprule
    \textbf{Prioritas} & \textbf{Kriteria} & \textbf{Contoh keputusan} \\
    \midrule
    P1 -- segera & Menghambat keputusan dan tenggat dekat & Lengkapi baseline, atasi prasyarat, tutup remedial tertunda \\
    P2 -- terjadwal & Berdampak besar tetapi memerlukan siklus latihan & Penguatan subtes, paket waktu, konsultasi pilihan \\
    P3 -- pengayaan & Fondasi relatif stabil dan tidak mendesak & Soal integratif, simulasi penuh, perluasan konteks \\
    \bottomrule
  \end{tabularx}}
  \vspace{0.14cm}
  {\color{StudentGray}\fontsize{9.7}{12}\selectfont Label P1--P3 adalah klasifikasi internal StudentPro, bukan tingkat resmi SNPMB.}
\end{frame}

\begin{frame}{StudentPro harus menghubungkan bukti dan tindakan}
  \kicker{Arsitektur Penutupan Gap}
  {\fontsize{10.1}{12.6}\selectfont
  \renewcommand{\arraystretch}{0.98}
  \begin{tabularx}{\textwidth}{>{\bfseries\color{StudentBlue}}p{0.19\textwidth} >{\RaggedRight\arraybackslash}X >{\RaggedRight\arraybackslash}p{0.29\textwidth}}
    \toprule
    \textbf{Fungsi LMS} & \textbf{Data yang dihubungkan} & \textbf{Keputusan} \\
    \midrule
    Profil siswa & Kohor, target, baseline, dan pilihan studi & Prioritas awal \\
    Pembelajaran & Materi, indikator, latihan, waktu, dan kesalahan & Remedial atau pengayaan \\
    Simulasi & Skor, ritme, soal kosong, dan stabilitas & Strategi ujian \\
    Konsultasi & Masalah, kesepakatan, tenggat, dan status & Pendampingan individual \\
    Pelaporan & Tren dan tindak lanjut terpilih & Komunikasi sekolah dan orang tua \\
    \bottomrule
  \end{tabularx}}
\end{frame}

\begin{frame}{Setiap kesenjangan memiliki kartu tindak lanjut}
  \kicker{Data Minimum}
  \begin{columns}[T,onlytextwidth]
    \begin{column}{0.46\textwidth}
      \begin{itemize}
        \item Siswa, subtes, domain, dan indikator.
        \item Bukti awal: tes, latihan, atau observasi.
        \item Jenis kesalahan dan tingkat prioritas.
        \item Tindakan serta penanggung jawab.
      \end{itemize}
    \end{column}
    \begin{column}{0.48\textwidth}
      \begin{itemize}
        \item Tenggat dan status pelaksanaan.
        \item Bukti tes ulang atau perubahan perilaku.
        \item Keputusan lanjut, ulang, atau pengayaan.
        \item Catatan komunikasi dengan siswa/orang tua.
      \end{itemize}
    \end{column}
  \end{columns}
  \vspace{0.20cm}
  \begin{beamercolorbox}[rounded=true,sep=0.18cm,wd=\textwidth]{block body}
    \centering\bfseries Satu gap memiliki satu pemilik, satu tenggat, dan satu bukti penutupan.
  \end{beamercolorbox}
\end{frame}

\begin{frame}{Temuan akademik menghasilkan keputusan berbeda}
  \kicker{Contoh Aturan Tindak Lanjut -- Bagian 1}
  {\fontsize{10.1}{12.6}\selectfont
  \renewcommand{\arraystretch}{0.99}
  \begin{tabularx}{\textwidth}{>{\bfseries\color{StudentBlue}\RaggedRight\arraybackslash}p{0.32\textwidth} >{\RaggedRight\arraybackslash}X}
    \toprule
    \textbf{Temuan} & \textbf{Keputusan awal} \\
    \midrule
    Konsep prasyarat salah berulang & Kembali ke materi dasar dan tes ulang topikal \\
    Akurasi baik, waktu berlebih & Latihan strategi dan batas waktu bertahap \\
    Hasil topikal baik, simulasi penuh turun & Latihan stamina dan ritme antarsubtes \\
    \bottomrule
  \end{tabularx}}
\end{frame}

\begin{frame}{Temuan proses menghasilkan keputusan berbeda}
  \kicker{Contoh Aturan Tindak Lanjut -- Bagian 2}
  {\fontsize{10.5}{13.0}\selectfont
  \renewcommand{\arraystretch}{1.02}
  \begin{tabularx}{\textwidth}{>{\bfseries\color{StudentBlue}\RaggedRight\arraybackslash}p{0.32\textwidth} >{\RaggedRight\arraybackslash}X}
    \toprule
    \textbf{Temuan} & \textbf{Keputusan awal} \\
    \midrule
    Aktivitas tidak tuntas & Identifikasi hambatan jadwal, perangkat, atau motivasi \\
    Pilihan studi berubah-ubah & Konsultasi berbasis minat, profil kemampuan, dan jalur \\
    \bottomrule
  \end{tabularx}}
  \vspace{0.24cm}
  \claim{Tindakan selalu mengikuti penyebab yang paling mungkin dan diperiksa kembali melalui bukti baru.}
\end{frame}

\begin{frame}{Penutupan kesenjangan memiliki pembagian peran}
  \kicker{Siapa Melakukan Apa?}
  {\fontsize{9.9}{12.3}\selectfont
  \renewcommand{\arraystretch}{0.97}
  \begin{tabularx}{\textwidth}{>{\bfseries\color{StudentBlue}}p{0.20\textwidth} >{\RaggedRight\arraybackslash}X}
    \toprule
    \textbf{Pihak} & \textbf{Tanggung jawab utama} \\
    \midrule
    Sekolah & Menetapkan kohor, jadwal, kebijakan, dan koordinasi tindak lanjut \\
    Tentor & Mendiagnosis kesalahan, memberi intervensi, dan memeriksa tes ulang \\
    Siswa & Menjalankan tugas, mencatat hambatan, dan merefleksikan hasil \\
    Orang tua & Menjaga rutinitas, kesehatan, ketenangan, dan komunikasi \\
    StudentPro & Menghubungkan data, notifikasi, status, dan laporan keputusan \\
    \bottomrule
  \end{tabularx}}
\end{frame}

\begin{frame}{Laporan orang tua berfokus pada arah perbaikan}
  \kicker{Komunikasi Tanpa Membandingkan Anak}
  \begin{columns}[T,onlytextwidth]
    \begin{column}{0.46\textwidth}
      {\color{StudentBlue}\bfseries Isi laporan}\par
      \begin{itemize}
        \item Kekuatan dan prioritas saat ini.
        \item Tren akurasi, waktu, dan ketuntasan.
        \item Tindak lanjut serta tenggat.
        \item Bukti perbaikan setelah tes ulang.
      \end{itemize}
    \end{column}
    \begin{column}{0.48\textwidth}
      {\color{StudentBlue}\bfseries Respons keluarga}\par
      \begin{itemize}
        \item Menanyakan hambatan secara tenang.
        \item Membantu konsistensi jadwal.
        \item Menghindari label ``pintar'' atau ``lemah''.
        \item Hadir pada keputusan pilihan studi.
      \end{itemize}
    \end{column}
  \end{columns}
  \vspace{0.15cm}
  \claim{Orang tua menerima apa yang perlu dibantu, bukan daftar kekurangan anak.}
\end{frame}

\begin{frame}{Penutupan kesenjangan dilakukan bertahap}
  \kicker{Urutan Implementasi}
  \begin{columns}[T,onlytextwidth]
    \begin{column}{0.22\textwidth}
      \begin{block}{1. Rapikan}
        Finalisasi kohor, identitas, sumber data, dan status validasi.
      \end{block}
    \end{column}
    \begin{column}{0.22\textwidth}
      \begin{block}{2. Petakan}
        Lengkapi baseline, target, indikator, serta jenis kesalahan.
      \end{block}
    \end{column}
    \begin{column}{0.22\textwidth}
      \begin{block}{3. Tindak}
        Jalankan prioritas, remedial, strategi, dan konsultasi.
      \end{block}
    \end{column}
    \begin{column}{0.22\textwidth}
      \begin{block}{4. Buktikan}
        Tes ulang, tinjau hasil, dan tutup atau perbarui status gap.
      \end{block}
    \end{column}
  \end{columns}
  \vspace{0.23cm}
  \centering{\color{StudentNavy}\bfseries\fontsize{13.2}{16}\selectfont Urutan ini menjaga program bergerak dari data menuju bukti perubahan.}
\end{frame}

\begin{frame}{Keberhasilan diukur sebelum hasil seleksi keluar}
  \kicker{Indikator Penutupan Gap}
  {\fontsize{10.0}{12.5}\selectfont
  \renewcommand{\arraystretch}{0.99}
  \begin{tabularx}{\textwidth}{>{\bfseries\color{StudentBlue}}p{0.23\textwidth} >{\RaggedRight\arraybackslash}X}
    \toprule
    \textbf{Indikator} & \textbf{Bukti yang dicari} \\
    \midrule
    Cakupan data & Seluruh kohor memiliki profil, baseline, dan target yang valid \\
    Ketuntasan proses & Materi, latihan, serta tindak lanjut selesai tepat waktu \\
    Kemajuan akademik & Akurasi, waktu, dan stabilitas membaik pada tes setara \\
    Kecepatan respons & Jarak antara temuan, intervensi, dan tes ulang semakin pendek \\
    Kualitas keputusan & Program belajar dan pilihan studi dapat dijelaskan dengan bukti \\
    \bottomrule
  \end{tabularx}}
\end{frame}

\begin{frame}{Kesimpulan analisis kesenjangan}
  \kicker{Pesan Utama Bab 8}
  \begin{columns}[T,onlytextwidth]
    \begin{column}{0.30\textwidth}
      \begin{block}{Jujur}
        Keterbatasan data dinyatakan; angka tidak dilengkapi dengan perkiraan.
      \end{block}
    \end{column}
    \begin{column}{0.30\textwidth}
      \begin{block}{Spesifik}
        Gap diuraikan menurut subtes, proses, strategi, tindak lanjut, dan pilihan studi.
      \end{block}
    \end{column}
    \begin{column}{0.30\textwidth}
      \begin{block}{Dapat ditutup}
        Setiap prioritas memiliki tindakan, pemilik, tenggat, dan bukti pemeriksaan ulang.
      \end{block}
    \end{column}
  \end{columns}
  \vspace{0.27cm}
  \centering{\color{StudentNavy}\bfseries\fontsize{14}{17}\selectfont Bab 9 mengubah prioritas ini menjadi target program UTBK 2027.}
\end{frame}
